%% Chapter 6 — Firmware Architecture
\chapter{Firmware Architecture}
\label{ch:firmware}

\section{Firmware Structure}

The access node firmware is located in \texttt{v2/access\_node/}. It is organised as a collection of single-responsibility modules coordinated by a central state machine.

\begin{table}[H]
\centering
\caption{Firmware Module Map}
\label{tab:modules}
\begin{tabular}{L{4.5cm}L{4.5cm}L{5.5cm}}
\toprule
\textbf{Module} & \textbf{Files} & \textbf{Responsibility} \\
\midrule
Config & \texttt{config.h} & Machine ID, name, session limit, broker IP \\
Secrets & \texttt{secrets.h} & WiFi credentials, master key (gitignored) \\
NVS Manager & \texttt{nvs\_manager.h/.cpp} & Read/write NVS namespace \texttt{mms\_config} \\
RFID Handler & \texttt{rfid\_handler.h/.cpp} & MFRC522 SPI init; card detect; sector key derivation; block read \\
Card Schema & \texttt{card\_schema.h/.cpp} & Parse blocks 4+5; validate schema; return \texttt{CardData} struct \\
Access Control & \texttt{access\_control.h/.cpp} & Revocation list check + permission bitmask check \\
Relay Controller & \texttt{relay\_controller.h/.cpp} & GPIO abstraction; safe OFF on init \\
Slot Sensor & \texttt{slot\_sensor.h/.cpp} & HC89 debounced input \\
Session Manager & \texttt{session\_manager.h/.cpp} & Session start/end/timeout; build log entry JSON \\
MQTT Client & \texttt{mqtt\_client.h/.cpp} & PubSubClient wrapper; reconnect backoff; command callback \\
Revocation Cache & \texttt{revocation\_cache.h/.cpp} & LittleFS \texttt{/config/revoked.json}; UID lookup \\
LittleFS Logger & \texttt{littlefs\_logger.h/.cpp} & JSONL append; rotate at 50KB; flush on connect \\
Display Primary & \texttt{display\_primary.h/.cpp} & 128$\times$64 SPI OLED; state-driven rendering \\
Display Status & \texttt{display\_status.h/.cpp} & 128$\times$32 I2C OLED; persistent status strip \\
LED Controller & \texttt{led\_controller.h/.cpp} & WS2812; non-blocking animations \\
Watchdog & \texttt{watchdog.h/.cpp} & ESP32 TWDT; 30s timeout; feed in main loop \\
OTA Handler & \texttt{ota\_handler.h/.cpp} & Version check; \texttt{esp\_https\_ota}; rollback \\
State Machine & \texttt{state\_machine.h/.cpp} & Central FSM; drives all transitions \\
Main & \texttt{access\_node.ino} & \texttt{setup()} + \texttt{loop()}; watchdog feed \\
\bottomrule
\end{tabular}
\end{table}

\section{Firmware State Machine}

The state machine is the central coordinator. All hardware events and MQTT messages cause state transitions. No module changes hardware state directly --- all changes go through the state machine.

\begin{figure}[H]
\centering
\begin{tikzpicture}[node distance=8mm and 9mm]
\node[state] (boot) {BOOT};
\node[state, below=of boot] (idle) {IDLE};
\node[state, below=of idle] (read) {CARD\\READING};
\node[mmswarn, left=of read] (denied) {ACCESS\\DENIED};
\node[mmswarn, right=of read] (disabled) {MACHINE\\DISABLED};
\node[state, below=of read] (active) {SESSION\\ACTIVE};
\node[mmswarn, below left=of active] (warning) {TIMEOUT\\WARNING};
\node[mmswarn, below=of warning] (expired) {TIMEOUT\\EXPIRED};
\node[state, below=of active] (ended) {SESSION\\ENDED};
\node[state, right=of active] (ota) {OTA\\UPDATE};
\draw[mmsline] (boot) -- node[right,font=\scriptsize]{init hardware, WiFi, MQTT} (idle);
\draw[mmsline] (idle) -- node[right,font=\scriptsize]{HC89 card detected} (read);
\draw[mmsline] (read) -- node[above,font=\scriptsize]{auth/schema/permission fail} (denied);
\draw[mmsline] (read) -- node[above,font=\scriptsize]{machine inactive} (disabled);
\draw[mmsline] (read) -- node[right,font=\scriptsize]{permission OK, relay ON} (active);
\draw[mmsline] (active) -- node[left,font=\scriptsize]{limit - 2 min} (warning);
\draw[mmsline] (warning) -- node[left,font=\scriptsize]{limit reached} (expired);
\draw[mmsline] (expired) -| node[below,font=\scriptsize]{relay OFF} (ended);
\draw[mmsline] (active) -- node[right,font=\scriptsize]{card removed / stop} (ended);
\draw[mmsline] (ended.west) -- ++(-34mm,0) |- node[pos=0.15,below,font=\scriptsize]{after display timeout} (idle.west);
\draw[mmsdash] (denied) |- (idle);
\draw[mmsdash] (disabled) |- (idle);
\draw[mmsdash] (active) -- node[above,font=\scriptsize]{MQTT OTA} (ota);
\draw[mmsdash] (ota) |- (boot);
\end{tikzpicture}
\caption{Access Node Firmware State Machine}
\label{fig:state_machine}
\end{figure}

\section{NVS Configuration}

NVS (Non-Volatile Storage) persists configuration across reboots. The namespace is \texttt{mms\_config}.

\begin{table}[H]
\centering
\caption{NVS Configuration Keys}
\begin{tabular}{L{3.5cm}C{2.5cm}L{3cm}L{5cm}}
\toprule
\textbf{Key} & \textbf{Type} & \textbf{Set by} & \textbf{Description} \\
\midrule
\texttt{machine\_id}   & uint8   & \texttt{config.h} at flash & 1--32 \\
\texttt{master\_key}   & blob (32B) & \texttt{secrets.h} at flash & HMAC master secret \\
\texttt{machine\_name} & string (32) & MQTT config push & Display name; fallback to \texttt{config.h} \\
\texttt{session\_limit} & uint16 & MQTT config push & Minutes; fallback to \texttt{config.h} \\
\texttt{relay\_active\_high} & uint8 & MQTT config push & 1=HIGH activates relay; 0=LOW \\
\texttt{machine\_active} & uint8 & MQTT config push & 0=disabled; all access denied \\
\texttt{fw\_version}  & string & OTA handler & Current firmware version string \\
\bottomrule
\end{tabular}
\end{table}

On first boot (no NVS data), the firmware uses the fallback values from \texttt{config.h}. On the first MQTT connect, the bridge pushes the Firestore config, which overwrites NVS. On subsequent boots without MQTT connection, NVS values are used.

\section{LittleFS File Layout}

\begin{table}[H]
\centering
\caption{LittleFS file layout}
\begin{tabular}{L{5cm}L{9cm}}
\toprule
\textbf{Path} & \textbf{Purpose} \\
\midrule
\texttt{/logs/sessions.jsonl} & Append-only offline session event log. \\
\texttt{/logs/sessions.old.jsonl} & Previous rotation backup; overwritten by the next rotation. \\
\texttt{/config/revoked.json} & Cached revoked UID list with \texttt{updated\_at} timestamp. \\
\bottomrule
\end{tabular}
\end{table}

\textbf{Log rotation:} When \texttt{sessions.jsonl} exceeds 50KB:
\begin{enumerate}[leftmargin=1.5em]
    \item Rename \texttt{sessions.jsonl} $\rightarrow$ \texttt{sessions.old.jsonl} (overwrites previous)
    \item Create new empty \texttt{sessions.jsonl}
    \item Continue appending
\end{enumerate}

\textbf{Flush on MQTT connect:}
\begin{enumerate}[leftmargin=1.5em]
    \item Open \texttt{sessions.jsonl}
    \item Read line by line; attempt JSON parse each line
    \item Discard any line that fails to parse (partial write from power loss)
    \item Publish each valid line to \texttt{mms/nodes/\{id\}/event} (QoS 1)
    \item Wait for broker ACK before publishing next
    \item After all lines published: delete file
\end{enumerate}

\section{Session Log Entry Format}

Each line in \texttt{sessions.jsonl} is a compact JSON object:

\begin{lstlisting}[style=jsonlstyle,caption={Session log entry format}]
{"t":1718600000,"r":"220002123","m":3,"e":"session_end",
 "d":245,"er":"card_removed"}
\end{lstlisting}

Fields:
\begin{itemize}[leftmargin=1.5em,noitemsep]
    \item \texttt{t}: Unix timestamp of session end
    \item \texttt{r}: Roll number string
    \item \texttt{m}: Machine ID integer
    \item \texttt{e}: Event type (\texttt{"session\_end"}, \texttt{"boot"})
    \item \texttt{d}: Duration in seconds
    \item \texttt{er}: End reason (\texttt{"card\_removed"}, \texttt{"timeout"}, \texttt{"remote\_stop"}, \texttt{"watchdog"})
\end{itemize}

\section{Display Content}

\subsection{Primary OLED (128$\times$64)}

\begin{table}[H]
\centering
\caption{Primary Display Content by State}
\begin{tabular}{L{4cm}L{10.5cm}}
\toprule
\textbf{State} & \textbf{Display content} \\
\midrule
IDLE & ``Tinkerers' Lab / IIT Indore'' + machine name + ``Tap card to start'' \\
CARD\_READING & ``Reading card...'' + spinner animation \\
ACCESS\_DENIED & ``!! Access Denied !!'' + reason + roll number (clears after 3s) \\
MACHINE\_DISABLED & ``Machine Offline'' + ``Contact admin'' \\
SESSION\_ACTIVE & Roll number + machine name + ``Time: MM:SS / HH:MM'' \\
TIMEOUT\_WARNING & ``!! 2 min left !!'' + ``Remove card soon'' + roll number \\
SESSION\_ENDED & ``Session Complete'' + roll + duration + ``Thank you!'' \\
OTA\_UPDATE & ``Updating FW...'' + version + progress percentage + ``Do not remove'' \\
ERROR & ``System Error'' + error code + ``Restarting...'' \\
\bottomrule
\end{tabular}
\end{table}

\subsection{Status Strip OLED (128$\times$32)}

Always shows:
\begin{table}[H]
\centering
\caption{Status strip layout}
\begin{tabular}{L{4cm}L{10cm}}
\toprule
\textbf{Region} & \textbf{Content} \\
\midrule
Left indicators & WiFi and MQTT symbols; filled when connected and empty when disconnected. \\
Center text & Machine name from NVS or Firestore config push. \\
Right text & Current time after NTP sync. \\
\bottomrule
\end{tabular}
\end{table}

\section{WS2812 LED State Map}

\begin{table}[H]
\centering
\caption{LED Colour and Pattern per State}
\begin{tabular}{L{3.5cm}L{3cm}L{7cm}}
\toprule
\textbf{State} & \textbf{Colour} & \textbf{Pattern} \\
\midrule
IDLE & White (dim) & Slow breathe, 2s period \\
CARD\_READING & Cyan & Fast blink, 200ms \\
SESSION\_ACTIVE & Green & Solid \\
TIMEOUT\_WARNING & Yellow $\leftrightarrow$ Orange & Alternate, 1s \\
TIMEOUT\_EXPIRED & Red & Solid (500ms) \\
ACCESS\_DENIED & Red & 3$\times$ quick flash then off \\
MACHINE\_DISABLED & Purple & Solid dim \\
OTA\_UPDATE & Blue & Chasing animation \\
ERROR & Red & Solid \\
\bottomrule
\end{tabular}
\end{table}

\section{Watchdog}

The hardware watchdog is the last line of defence against firmware hangs:

\begin{lstlisting}[language=C,caption={Watchdog initialisation in setup()}]
esp_task_wdt_init(30, true);  // 30s timeout, panic on trigger
esp_task_wdt_add(NULL);       // watch the main task
\end{lstlisting}

\begin{lstlisting}[language=C,caption={Watchdog feed in main loop}]
void loop() {
    esp_task_wdt_reset();   // must be called at least every 30s
    stateMachine.tick();
    // ...
}
\end{lstlisting}

If \texttt{loop()} takes longer than 30 seconds (e.g., blocking I2C, infinite RFID retry), the TWDT triggers a panic and the ESP32 reboots. The boot reason is available via \texttt{esp\_reset\_reason()} on the next boot.
